\section{From collar roots to the completed universal tower}
\label{sec:realization}

The first stage, Proposition~\ref{prop:first-stage}, places the field
$L_\ell$ inside $H$.  We now use the auxiliary covers from
Proposition~\ref{prop:arbitrary-roots} to realize every finite Galois
extension of $K_0$ in that same completion.  The covers used to create
the collars may be branched; their horizontal ramification indices
divide $\ell$, so the first stage is sufficient to accommodate them.

\begin{proof}[Proof of Theorem~\ref{thm:completion}]
Let $F/K_0$ be any finite Galois extension inside $K^{\mathrm{sep}}$,
and let $C\to X_{\bar k}$ be its smooth proper curve cover.  Its branch
set $B$ is finite.  If $B$ is empty, then $F\subset\widetilde K\subset
H$ by definition of the universal étale field.  Suppose henceforth
that $B$ is nonempty.

For each $b\in B$, choose a rational local parameter $t_b$ at $b$ and
disjoint disks $B_b$ as in Lemma~\ref{lem:local-normal-form}; these
disks avoid $\xi$.  Let $m_b$ be the least common multiple of the
ramification indices of $C$ at the points above $b$.  Choose
$\beta_b\in\bar k^\times$ whose absolute value is smaller than the
radius of $B_b$ in the coordinate $t_b$.

Apply Proposition~\ref{prop:arbitrary-roots} with $m=m_b$.  In its
notation choose a positive integer $D_b>v_p(m_b)+1/(p-1)$ and
$\alpha_b\in\bar k^\times$ with
$a_b=-\log_p|\alpha_b|>2D_b$.  It gives a smooth proper connected
curve cover of the $T$-line whose horizontal indices divide $\ell$
and on whose full inverse image of
\[
                     D_b<-\log_p|T|<a_b-D_b
\]
the coordinate $T$ has an analytic $m_b$-th root.  Pull this cover back
by the nonconstant morphism $X_{\bar k}\to\mathbb P^1_{\bar k}$
given by $T=t_b/\beta_b$, and take a normalized connected component.
Denote the resulting cover by $W_b\to X_{\bar k}$.

Lemma~\ref{lem:horizontal-closure} shows that its horizontal indices
still divide $\ell$.  Roots pull back to every point of the normalized
component, so $t_b$ has an analytic $m_b$-th root on the entire inverse
image in $W_{b,\Omega}^{\mathrm{an}}$ of the collar
\[
 |\beta_b|p^{-(a_b-D_b)}<|t_b|<|\beta_b|p^{-D_b}
\]
inside $B_b$.  Here we multiplied the pulled-back root of $T$ by an
$m_b$-th root of $\beta_b$ in $\Omega$.

Embed the finitely many function fields $\bar k(W_b)$ into the fixed
$K^{\mathrm{sep}}$ over $K_0$ and form their compositum.  Let
$W\to X_{\bar k}$ be the associated smooth proper connected curve.
It dominates every $W_b$; its horizontal indices divide $\ell$ by
Lemma~\ref{lem:horizontal-closure}.  All collar roots therefore pull
back to the full inverse images in $W_\Omega^{\mathrm{an}}$.

Proposition~\ref{prop:patching}, with $Y=W_\Omega$, produces a
connected finite étale cover $V_\Omega\to W_\Omega$ and an analytic
map from its exterior region to $C_\Omega^{\mathrm{an}}$ over
$X_\Omega^{\mathrm{an}}$.  Since $W$ is proper over the algebraically
closed field $\bar k$, invariance of finite étale covers under the
algebraically closed extension $\bar k\subset\Omega$ \cite{SGA1}
gives a connected finite étale cover $V\to W$ and an isomorphism
\[
                  V\times_{\bar k}\Omega\simeq V_\Omega
                  \quad\text{over }W_\Omega.
\]
No descent assertion is needed for the exterior analytic map.

Although no such bound was imposed on $C\to X_{\bar k}$, the composite
$V\to X_{\bar k}$ still has horizontal indices dividing $\ell$.
By Lemma~\ref{lem:horizontal-closure}, its function
field embeds over $K_0$ in $L_\ell$.  This embedding can, if desired,
be chosen to extend the already chosen embedding of $\bar k(W)$:
extend that embedding to $\bar k(V)$ inside $K^{\mathrm{sep}}$ and
then take its Galois closure, which still has indices dividing $\ell$.
Proposition~\ref{prop:first-stage} places its image in $H$.

Lemma~\ref{lem:point-transfer} now gives a $K_0$-embedding
$F=\bar k(C)\hookrightarrow H\subset E$.  Because $F/K_0$ is
normal, this embedding has image the chosen subfield $F$ of $E$.
Indeed a primitive element of $F$ is sent to a root of its minimal
polynomial, all of whose roots already belong to $F$; the same holds
for every element.  The image, having degree $[F:K_0]$, is therefore
$F$ itself.  Consequently $F\subset H$.

Every element of $K^{\mathrm{sep}}$ lies in a finite Galois extension
of $K_0$ inside $K^{\mathrm{sep}}$: this is an algebraic closure of
the characteristic-zero field $K_0$.  Taking the union over those
extensions proves $K^{\mathrm{sep}}\subset H$ inside $E$.
\end{proof}
